\begin{theorem}
Let \(C>1\) be fixed. Then, for all sufficiently large \(s\),
\[
  r(s,\lceil Cs\rceil)\ge \left(2^{1-\frac{1}{2C}}\right)^s.
\]
\end{theorem}
\begin{proof}
Fix \(C>1\). For each \(s\), set
\[
  k=\lceil Cs\rceil.
\]
We will prove the desired inequality for all sufficiently large \(s\).

Since \(C>1\), after increasing the lower threshold for \(s\) if necessary, we have \(4\le s\le k\). By \(\noderef{F2ForwardIndependentTuples}\), there are a finite vertex set \(W\) and a digraph \(D\) on \(W\) such that \(D\) is loopless, \(D\) is \(T_s\)-free, and
\[
  N\coloneqq |W|
  =2^{2s-3}-2^{s-1}-2^{s-2}+1
\]
for this same value of \(s\). For this same digraph \(D\) and integer \(k\), \(\noderef{F2ForwardIndependentLinearBound}\) gives, after perhaps increasing the lower threshold for \(s\) again, the additional estimate, in the notation of \(\noderef{ForwardIndependentTupleCount}\),
\[
  \overrightarrow{i_k}(D)
  \le 2^{sk+s^2/2}. \tag{1}
\]
For \(s\ge4\),
\[
  N=2\cdot 2^{2s-4}-3\cdot 2^{s-2}+1
  \ge 2^{2s-4}, \tag{2}
\]
so \(N=(1+o(1))2^{2s-3}\) in the coarse sense used below. Also, since \(C>1\), for all sufficiently large \(s\) we have \(k=\lceil Cs\rceil\ge1\).

Apply \(\noderef{DigraphToGraphIndependentSetBound}\) to \(D\), using the \(T_s\)-freeness from \(\noderef{F2ForwardIndependentTuples}\) and \(k\ge1\). We obtain a \(K_s\)-free simple graph \(G\) on the same vertex set \(W\) such that, writing \(i_k(G)\) for the independent-set count from \(\noderef{SimpleGraphIndependentSetCount}\),
\[
  i_k(G)\le \left(\frac e k\right)^k \overrightarrow{i_k}(D).
\]
Together with (1), this gives
\[
  i_k(G)\le \left(\frac e k\right)^k 2^{sk+s^2/2}. \tag{3}
\]

First suppose \(i_k(G)=0\). Then \(G\) is a graph on \(N\) vertices with no \(K_s\) and no independent set of size \(k\). By \(\noderef{RamseyProperty}\), this means the Ramsey property \(R(s,k,N)\) fails. Moreover, for every \(m\le N\), restricting \(G\) to any \(m\)-vertex subset still gives a graph with no \(K_s\) and no independent \(k\)-set, so \(R(s,k,m)\) fails as well. Since \(\noderef{RamseyNumber}\) defines \(r(s,k)\) as the least \(m\) for which \(R(s,k,m)\) holds, it follows that
\[
  r(s,k)>N.
\]
By (2), \(N\ge2^{2s-4}\), while \(1-\frac1{2C}<1\). Hence, for all sufficiently large \(s\),
\[
  N\ge 2^{(1-\frac1{2C})s}.
\]
The required lower bound follows in this case.

It remains to consider the case \(i_k(G)>0\). Define
\[
  p=i_k(G)^{-1/k}.
\]
Because \(i_k(G)\) is a positive integer, \(0<p\le1\), and by definition
\[
  p^k i_k(G)=1.
\]
Thus the hypotheses of \(\noderef{SamplingKsFreeRamseyBound}\) are satisfied for the \(K_s\)-free graph \(G\), the integer \(k\ge1\), and this value of \(p\). Therefore
\[
  r(s,k)>pN-1. \tag{4}
\]

We now lower-bound \(pN\). From (3),
\[
  p=i_k(G)^{-1/k}
  \ge
  \left(\left(\frac e k\right)^k2^{sk+s^2/2}\right)^{-1/k}
  =
  \frac{k}{e}\,2^{-(sk+s^2/2)/k}. \tag{5}
\]
Since \(k=\lceil Cs\rceil\) and \(C>1\), we have \(k\ge s\). Combining (2), (4), and (5) gives
\[
  r(s,k)>
  2^{2s-4}\frac{s}{e}
  2^{-(sk+s^2/2)/k}-1.
\]
The exponent simplifies as
\[
  2s-4-\frac{sk+s^2/2}{k}
  =
  s-4-\frac{s^2}{2k}.
\]
Since \(k\ge Cs\),
\[
  \frac{s^2}{2k}\le \frac{s}{2C}.
\]
Consequently
\[
  r(s,k)>
  \frac{s}{16e}\,
  2^{(1-\frac1{2C})s}-1. \tag{6}
\]
The factor \(s/(16e)\) tends to infinity. Therefore, after increasing the lower threshold for \(s\) one last time, the right side of (6) is at least
\[
  2^{(1-\frac1{2C})s}.
\]
This proves
\[
  r(s,\lceil Cs\rceil)\ge
  \left(2^{1-\frac1{2C}}\right)^s
\]
for all sufficiently large \(s\), as claimed.
\end{proof}
